\documentclass[11pt]{article}
\usepackage[margin=1in,top=0.85in,bottom=0.8in]{geometry}
\usepackage{amsmath,amssymb}
\usepackage[protrusion=true,expansion=false]{microtype}
\usepackage{xcolor}
\usepackage{enumitem}
\definecolor{acc}{HTML}{1B6B58}
\definecolor{grey}{HTML}{5C6763}
\newcommand{\E}{\mathbb{E}}
\newcommand{\Prob}{\mathbb{P}}
\newcommand{\Var}{\operatorname{Var}}
\newcommand{\Poi}{\mathrm{Poisson}}
\newcommand{\Ber}{\mathrm{Bernoulli}}
\newcommand{\TV}{d_{\mathrm{TV}}}
\pagestyle{empty}
\setlength{\parindent}{0pt}
\setlength{\parskip}{0.35em}

\begin{document}

{\footnotesize\color{acc}\textsc{ORF 526 \quad$\cdot$\quad Quiz 2 \quad$\cdot$\quad instructor key --- not for distribution}}

\vspace{0.4em}
{\large\bfseries Quiz 2 --- answers and marking notes}

\vspace{0.4em}
\hrule height 0.5pt
\vspace{0.4em}

\textbf{1.} \ (a) $\varphi_B(t)=(1-p)+pe^{it}=1+p(e^{it}-1)$.

(b) $\displaystyle\varphi_{N_n}(t)=\Big[1+\tfrac{\lambda}{n}(e^{it}-1)\Big]^{n}\longrightarrow \exp\big(\lambda(e^{it}-1)\big)$, the characteristic function of $\Poi(\lambda)$.

(c) \textbf{Independence}, which is what gives $\varphi_{N_n}=\varphi_{B}^{\,n}$.

{\color{grey}\emph{Marking.} Part (c) is the one that matters and the one most likely to be skipped; the whole point of the week is that this step is where independence is spent. Accept ``the $B_{n,i}$ are i.i.d.'' but probe whether the student knows that \emph{identical} distribution only lets them write one factor $n$ times --- it is independence that lets them multiply at all. Do not require the $(1+z/n)^n\to e^z$ lemma to be quoted by name.}

\vspace{0.9em}

\textbf{2.} \ (a) $\E W=\sum_i\Prob(\pi(i)=i)=n\cdot\tfrac1n=1$.

(b) \textbf{Independence of the indicators.} The $I_i$ are dependent: $\Prob(I_1=I_2=1)=\frac{1}{n(n-1)}$, whereas $\Prob(I_1=1)\Prob(I_2=1)=\frac{1}{n^2}$. Le Cam's hypothesis is not satisfied, so the theorem says nothing here.

(c) A theorem whose hypotheses fail proves nothing, whatever the truth of its conclusion. The bound has to be obtained another way --- in Homework 4, by a coupling.

{\color{grey}\emph{Marking.} This is the question that separates students who memorised a formula from students who read its hypotheses. Full credit in (b) requires naming \emph{independence} specifically, not ``the $I_i$ aren't nice'' or ``$W$ isn't a sum of Bernoullis'' --- it is a sum of Bernoullis, which is exactly the trap. In (c), the sentence must locate the flaw in the \emph{argument}, not in the conclusion; ``the bound is loose'' misses the point and earns nothing.}

\vspace{0.9em}

\textbf{3.} \ (a) $\displaystyle \E[W'-W\mid X_1,\dots,X_n]=-\tfrac{1}{\sqrt n}\,\E[X_I-X_I'\mid X_1,\dots,X_n]=-\tfrac{1}{\sqrt n}\Big(\tfrac1n\sum_iX_i-0\Big)=-\frac{W}{n}.$

(b) $a=1/n$, by the tower property.

(c) Because the bound of Homework 4 needs only $\E[W'-W\mid W]$ and $\E[(W'-W)^2\mid W]$ --- conditional moments of a coupling, which can be computed without ever factorising $\E[X_if(W^{(i)})]$ into $\E[X_i]\,\E[f(W^{(i)})]$.

{\color{grey}\emph{Marking.} In (a) the two things to check are that $\E[X_I'\mid X_1,\dots,X_n]=0$ (the fresh copy is independent of everything conditioned on) and that averaging over $I$ produces $\frac1n\sum_iX_i$. A student who writes $\E[X_I\mid X_1,\dots,X_n]=X_I$ has not seen that $I$ is random and should lose (a) but may still earn (b) and (c). Accept any (c) that identifies the payoff as ``only conditional moments of the pair are needed''.}

\vspace{1.1em}
\hrule height 0.4pt
\vspace{0.3em}
{\footnotesize\color{grey} Sources: Q1 from Homework 3, Problems 3 and 5 (core); Q2 from Homework 4, Problem 10 (core); Q3 from Homework 4, Problem 5 (core). Question types rotate against Quiz 1, which used \textsc{compute}, \textsc{compute}, \textsc{which hypothesis fails}.}

\end{document}
